% =============================================================================
% idea2/emi.tex — Perfil de Proyecto de Grado
% Estructura según: Guía de Elaboración de Perfil Trabajos de Grado
% Universidad Privada del Valle — DAAP, 1ra Ed. 2019
%
% Motor requerido: XeLaTeX
%   Compilar con: latexmk -xelatex emi.tex  (ejecutar desde esta carpeta)
% =============================================================================

\documentclass{univalle-perfil}

\graphicspath{{../}}
\addbibresource{emi.bib}

% -----------------------------------------------------------------------------
% METADATOS DE LA CARÁTULA
% -----------------------------------------------------------------------------
\universidad{Universidad Privada del Valle}
\facultad{Facultad de Informática y Electrónica}
\carrera{Carrera de Licenciatura en Ingeniería de Sistemas}
\tituloTrabajo{TODO: Título del proyecto}
\modalidad{Proyecto de Grado}
\tituloOptar{Licenciatura en Ingeniería de Sistemas}
\postulante{Franco Javier Torrez Alvarado}
\tutor{TODO: Nombre del tutor}
\ciudad{Santa Cruz}
\pais{Bolivia}
\anio{2026}

% =============================================================================
\begin{document}

% --- Carátula ----------------------------------------------------------------
\makecaratula

% --- Índices -----------------------------------------------------------------
\pagenumbering{gobble}
\tableofcontents
\listoffigures

% --- Contenido principal -----------------------------------------------------
\pagenumbering{arabic}

% =============================================================================
% 1. INTRODUCCIÓN
% =============================================================================
\section{Introducción}

% TODO: ~1 página. Fundamentación científica resumida. Novedad y actualidad
% del tema. Último párrafo: presentación del trabajo a realizar.

% =============================================================================
% 2. ANTECEDENTES DEL PROBLEMA
% =============================================================================
\section{Antecedentes del Problema}

% TODO: ~1 página. Contexto del problema. Proyectos similares existentes.
% Incluir diagrama Ishikawa y análisis FODA.

\subsection{Planteamiento del Problema}

% TODO: 1 párrafo de 3 a 5 líneas. Formulación del problema (puede ser
% en forma de pregunta). Especificar la población que se investiga.

\subsection{Objeto de Estudio}

% TODO: 1 a 2 líneas. Responde al "¿qué voy a estudiar?".

% =============================================================================
% 3. OBJETIVOS
% =============================================================================
\section{Objetivos}

\subsection{Objetivo General}

% TODO: Qué / Para qué / Cómo / Dónde. Debe ser medible.

\subsection{Objetivos Específicos}

\begin{itemize}
    \item % TODO
    \item % TODO
    \item % TODO
    \item % TODO
\end{itemize}

% =============================================================================
% 4. JUSTIFICACIÓN
% =============================================================================
\section{Justificación}

% TODO: ~1 página total distribuida entre las tres subsecciones.

\subsection{Justificación Técnica}

% TODO

\subsection{Justificación Económica}

% TODO

\subsection{Justificación Social}

% TODO

% =============================================================================
% 5. DELIMITACIÓN DE LA INVESTIGACIÓN
% =============================================================================
\section{Delimitación de la Investigación}

\subsection{Temática}

% TODO: Límite del estudio en cuanto a teorías, áreas del conocimiento,
% normas y regulaciones aplicables.

\subsection{Espacial}

% TODO: Lugar geográfico del estudio y lugar de aplicación de resultados.

\subsection{Temporal}

% TODO: Período de tiempo, fechas estimadas, límite de validez.

% =============================================================================
% 6. PROPUESTA DE SOLUCIÓN AL PLANTEAMIENTO DEL PROBLEMA
% =============================================================================
\section{Propuesta de Solución al Planteamiento del Problema}

% TODO: En forma gráfica exponer la posible solución. 1 a 2 párrafos de
% explicación. Incluir diagrama de módulos/arquitectura y diagrama de flujo.

% =============================================================================
% 7. DISEÑO METODOLÓGICO
% =============================================================================
\section{Diseño Metodológico}

\subsection{Tipo de Investigación}

% TODO: Identificar el/los tipos de investigación y explicar brevemente
% las definiciones teóricas aplicadas al trabajo y contexto.

\subsection{Método de Investigación}

% TODO: Métodos seleccionados (histórico-lógico, análisis documental,
% enfoque de sistemas, modelación, etc.) y para qué sirve cada uno.

\subsection{Técnicas e Instrumentos de Investigación}

\subsubsection{Técnicas}

% TODO: Qué técnica, a quién se aplica, para qué se utilizará.

\subsubsection{Instrumentos}

% TODO: Solo mencionar los instrumentos correspondientes a cada técnica.

\subsection{Población y Muestra}

% TODO: Población y muestra para cada técnica. Tipo de muestreo.

% =============================================================================
% 8. ÍNDICE TENTATIVO
% =============================================================================
\section{Índice Tentativo}

\begin{indicetentativo}
INTRODUCCIÓN \\[2pt]
CAPÍTULO I. MARCO GENERAL \\
\quad 1.1 Antecedentes \\
\quad 1.2 Planteamiento del Problema \\
\quad\quad 1.2.1 Formulación del Problema \\
\quad\quad 1.2.2 Árbol del Problema \\
\quad 1.3 Objetivos \\
\quad\quad 1.3.1 Objetivo General \\
\quad\quad 1.3.2 Objetivos Específicos \\
\quad 1.4 Justificación \\
\quad\quad 1.4.1 Justificación Técnica \\
\quad\quad 1.4.2 Justificación Económica \\
\quad\quad 1.4.3 Justificación Social \\
\quad 1.5 Metodología \\
\quad\quad 1.5.1 Enfoque de la Investigación \\
\quad\quad 1.5.2 Tipo de Investigación \\
\quad\quad 1.5.3 Método de Investigación \\
\quad\quad 1.5.4 Técnicas e Instrumentos de Investigación \\
\quad\quad 1.5.5 Fuentes de Información \\
\quad\quad 1.5.6 Propuesta \\[2pt]
CAPÍTULO II. MARCO TEÓRICO Y TECNOLÓGICO \\
\quad 2.1 TODO \\
\quad 2.2 TODO \\[2pt]
CAPÍTULO III. ANÁLISIS Y DISEÑO DEL SISTEMA PROPUESTO \\
\quad 3.1 TODO \\
\quad 3.2 TODO \\[2pt]
CAPÍTULO IV. CONSTRUCCIÓN E IMPLEMENTACIÓN DEL SISTEMA \\
\quad 4.1 TODO \\
\quad 4.2 TODO \\[2pt]
CONCLUSIONES \\
RECOMENDACIONES \\
REFERENCIAS BIBLIOGRÁFICAS \\
ANEXOS
\end{indicetentativo}

% =============================================================================
% 9. MARCO TEÓRICO
% =============================================================================
\section{Marco Teórico}

% TODO: Iniciar con un índice tentativo de los tópicos teóricos. Desarrollar
% algunos aspectos con posiciones teóricas y referencias bibliográficas.

\subsection{TODO: Primera área teórica}

% TODO

\subsection{TODO: Segunda área teórica}

% TODO

\subsection{TODO: Tercera área teórica}

% TODO

\subsection{TODO: Cuarta área teórica}

% TODO

% =============================================================================
% 10. CRONOGRAMA
% =============================================================================
\section{Cronograma}

% TODO: Tabla o diagrama de Gantt con etapas y fechas del proyecto.

% =============================================================================
% 11. REFERENCIAS BIBLIOGRÁFICAS
% =============================================================================
\referencias

% =============================================================================
% ANEXOS
% =============================================================================
\seccioncuerpo{ANEXOS}

\end{document}
